\chapter{Introduction}
\label{ch:introduction}

\begin{quote}\itshape
This chapter in brief: jet engines degrade invisibly, airlines monitor them through a handful
of measurements, and two schools of methods compete to interpret those measurements. This
project builds a level playing field to compare them. If you are new to the field, skim the
hypotheses table on first read and return to it after \cref{ch:background}.
\end{quote}

\keyidea{The field compares AI against weakly-implemented traditional baselines. The only honest comparison gives both families \emph{identical} data, metrics, splits and tuning budget; everything in this document follows from insisting on that symmetry.}


\section{Motivation}

A commercial jet engine degrades from its very first flight. Dust and pollution stick to
compressor blades, hot-section parts creep and oxidize, tip clearances open up, seals wear.
The airline cannot see any of this directly ---the engine stays on the wing for years--- so it
watches the engine's \emph{symptoms} instead: the temperatures, pressures, shaft speeds and fuel
flow recorded on every flight. The discipline that turns those symptoms into maintenance
decisions is called \emph{Engine Health Monitoring} (EHM). The stakes are large: a single shop
visit for an engine overhaul costs millions of dollars, and an unscheduled engine removal
disrupts an entire fleet schedule.

The classical analytical core of EHM, in use since the 1970s, is \emph{Gas Path Analysis}
(GPA)~\cite{Urban1975}. The idea: each physical fault (a dirty compressor, an eroded turbine)
leaves a characteristic fingerprint on the measured parameters; by inverting a thermodynamic
model of the engine, one can estimate \emph{which} component has deteriorated and \emph{by how
much} from the measurements alone. Around GPA the industry built a mature toolbox ---baseline
models, exponential trend smoothing, Kalman filters, expert fault-isolation rules%
~\cite{Volponi2014}--- which this project collectively calls \emph{traditional EHM}.

Over the last decade, machine learning (ML) has produced striking results in the adjacent research
field of prognostics: predicting the \emph{Remaining Useful Life} (RUL) of an engine, i.e.\ how
many flights remain before it must come off the wing, mostly on synthetic benchmark datasets
such as NASA's C-MAPSS~\cite{Saxena2008} and N-CMAPSS~\cite{AriasChao2021}. Yet the literature
has a structural weakness: \textbf{AI methods are almost never compared against a traditional
pipeline implemented with equal care}, on the same data, with the same metrics and the same
tuning effort. The result is a stream of comparisons against straw men, which makes it
impossible to know what AI actually contributes, under which conditions, and at what cost.

\section{Objective}

This project builds a reproducible testbed in which both EHM families ---traditional and
AI-based--- observe \emph{exactly the same data} from a synthetic fleet of CFM56-7B26 turbofans
and are scored with \emph{exactly the same metrics}. The methodological key is that the
synthetic data carries \textbf{complete ground truth}: the true health state of every engine
component (its efficiency and flow-capacity deviations, defined in \cref{sec:gpa}) at every
flight. With real airline data this is impossible ---nobody knows the true internal state of an
engine on the wing--- but with synthetic data one can measure not only whether a method detects
or predicts well, but whether it \emph{attributes the damage to the correct component}. That
attribution question is precisely where traditional GPA is weakest (the \emph{smearing}
phenomenon, \cref{sec:gpa}) and where AI is least transparent (the black-box problem).

\paragraph{Audience and register.} This document is a research work aimed at two readers at
once: the engineering organization deciding whether and how to integrate AI into its engine
health monitoring practice, and the researcher looking for a rigorous, reproducible
benchmark at the knowledge frontier. That dual audience sets the register: every method is
pushed to publishable rigor (pre-registration, statistical tests, audited datasets), and
every concept is introduced without assuming prior specialization --- chapter briefs, worked
numeric examples and a notation front-matter keep the document self-contained. Advancing the
frontier and being readable are treated as the same job, not competing ones.

Since neither a performance model of the engine nor degradation data were available to the
project, both are built from scratch: a thermodynamic ``digital twin'' of the CFM56-7B26 in
pyCycle~\cite{Hendricks2019}, calibrated against certified public data (\cref{ch:model}), and a
fleet degradation generator with physically derived fault signatures (\cref{ch:fleet}).

\section{Hypotheses}
\label{sec:hypotheses}

Each hypothesis is formalized with quantitative thresholds \emph{before} any test-set result is
seen (the pre-registration protocol of \cref{sec:safeguards}). All are tested with a Wilcoxon
signed-rank test paired by engine, Holm correction for multiple comparisons
($\alpha = 0.05$), and BCa bootstrap confidence intervals. In plain words: for every comparison
we check that the improvement is statistically significant, not a fluke of one lucky engine.
The table below necessarily uses terms defined later ---the GPA quantities in
\cref{ch:background}, the machine-learning methods in \cref{sec:ai-toolbox}, every metric in
the metrics section of \cref{ch:methodology}, and all acronyms in the notation chapter--- so
it is meant to be skimmed now and revisited after those stops.

\begin{table}[htbp]
  \centering\small
  \caption{Project hypotheses and their confirmation criteria.}
  \label{tab:hypotheses}
  \begin{tabular}{p{0.04\textwidth}p{0.36\textwidth}p{0.50\textwidth}}
    \toprule
    ID & Hypothesis & Confirmation criterion \\
    \midrule
    H1 & AI detects faults earlier and with fewer false alarms than trend monitoring &
      median lead time $\geq +20$\,\% and false-positive rate $\leq 0.5\times$ versus
      trending, at a fixed recall of $0.9$; $p<0.05$ \\
    H2 & AI isolates faults with overlapping signatures better &
      $+10$ percentage points of isolation accuracy on the ``confusable'' subset (fault
      signatures less than $15^{\circ}$ apart in the influence-matrix space) versus
      weighted-least-squares GPA; $p<0.05$ \\
    H3 & AI predicts RUL better &
      simultaneous improvement in RMSE and in the asymmetric NASA PHM08 score; $p<0.05$ \\
    H4 & The physics-informed hybrid dominates when data is scarce &
      RUL RMSE $\leq$ pure AI with 100\,\% of the training data and strictly better with
      10\,\% and 25\,\%; $p<0.05$ \\
    H5 & Conformal intervals are better calibrated than the Kalman covariance &
      smaller $|$empirical coverage $-$ 0.9$|$ with mean interval width no worse than
      $+20$\,\% \\
    \bottomrule
  \end{tabular}
\end{table}

Refuting any hypothesis is itself a publishable result: pre-registration is what makes the
refutation credible. Reformulating hypotheses after the repository tag \code{prereg-v1} is
forbidden. One operationalization was refined \emph{at} the freeze, before any test-set run
and disclosed in the frozen document: H1's ``fixed recall $0.9$'' proved unreachable at the
v1.1 fault magnitudes, so its operating point was redefined per family under a false-alarm
budget (\cref{sec:res-h1}). Changing an operating definition at freeze time, on the record,
is legitimate; changing a threshold after seeing the answer is what the tag forbids.

\section{Contributions}
\label{sec:contributions}

In plain terms first, before the technical names: the project builds a realistic test dataset
with hidden answers no real fleet can provide (C1), runs a scrupulously fair contest between
the two method families (C2), and from it delivers a physics-informed hybrid (C3), calibrated
uncertainty (C4), a physics-anchored explainability score (C5), a sensor-and-data trade map
(C6), an external reality check (C7), and --- the furthest reach --- two per-engine
certificates, proven against the truth, of what a diagnosis can honestly know (C8) and how
much of an engine's future its present can reveal (C9). Each item below names
the technique; every term is defined where it is first used, so this list is a preview, not a
prerequisite.

\begin{enumerate}
  \item[\textbf{C1.}] \textbf{SynCFM56}: the first public synthetic GPA dataset for a specific
    commercial engine (CFM56-7B26), with per-flight health-parameter ground truth and an
    ACARS-style snapshot format (one takeoff report and one cruise report per flight, the way
    real airline EHM data actually arrives), more realistic than the continuous time series of
    C-MAPSS.
  \item[\textbf{C2.}] A head-to-head comparison under a \textbf{pre-registered protocol} with a
    symmetric tuning budget (the same number of Optuna trials) for both families.
  \item[\textbf{C3.}] A physics-informed hybrid with three separately ablated mechanisms:
    residual features against the digital twin, physically constrained loss functions, and
    GPA$\rightarrow$ML stacking.
  \item[\textbf{C4.}] Compared uncertainty quantification: conformal-prediction RUL intervals
    versus the Kalman filter covariance.
  \item[\textbf{C5.}] The \textbf{Physics-Consistency Score (PCS)}: a new explainability metric
    that projects SHAP attributions into the physical health-parameter space through the
    influence coefficient matrix.
  \item[\textbf{C6.}] A determinability map: a sensors $\times$ noise $\times$ data-size
    ablation showing \emph{when} each approach wins --- the sensor axis throughout
    (\cref{sec:icm,ch:gpa-practice}), the data-size axis in the H4 ablation
    (\cref{sec:ai-hybrid}), and the noise axis swept from half to four times nominal sensor
    quality in \cref{sec:noise-sweep}.
  \item[\textbf{C7.}] Sim-to-real validation: the full methodology is re-run on N-CMAPSS to
    check that the method ranking is not an artifact of our own simulator.
  \item[\textbf{C8.}] \textbf{The earned-confidence certificate} (\cref{ch:breakthrough}): a
    per-engine, history-resolved identifiability guarantee for gas-path diagnosis --- and,
    uniquely, a \emph{proof against ground truth} that it is honest ($\rho = 0.70$ between
    certified precision and actual error). It states exactly which health directions an
    engine's data lets one conclude, and doubles as a sensor-acquisition instrument
    (predicting that station probes rescue HPC efficiency $45\times$). The validation is one
    only a synthetic-truth benchmark can perform.
  \item[\textbf{C9.}] \textbf{The prognostic floor} (\cref{sec:bt-prognostic-floor}): the
    predictive counterpart of C8. By conditioning on each engine's \emph{true} current health
    state it measures, against ground truth, the aleatoric floor on remaining-life error --- the
    RMSE no method can beat. It shows $87\,\%$ of mid-life RUL uncertainty is irreducible, yet a
    $4\times$ reducible gap opens late in life, telling an operator \emph{where} a better
    prognostic is worth building and where the future is simply unknowable.
\end{enumerate}

\section{Document structure}

The document follows the work in order, and a reader can stop wherever their interest ends.
\Cref{ch:background} builds the foundations from zero: how a turbofan works, what is measured
on the wing, and the mathematics of gas-path diagnosis --- no prior knowledge assumed.
\Cref{ch:methodology} lays out the testbed and the fairness safeguards that make the
comparison trustworthy. The next chapters build the ingredients: the thermodynamic twin
of the engine (\cref{ch:model}), the classical gas-path diagnosis measured end to end on that
twin --- the quantitative floor everything later is compared against
(\cref{ch:gpa-practice}) --- what the data \emph{is} and where it comes from
(\cref{ch:data}), and the synthetic fleet generated from it (\cref{ch:fleet}). The two EHM
families are then implemented and measured --- traditional (\cref{ch:trad}) and AI
(\cref{ch:ai}) --- and adjudicated head to head under the pre-registered protocol
(\cref{ch:results}), with five real engines telling the story in operations language
(\cref{ch:cases}).

The remaining chapters push past the main study. \Cref{ch:tomography} asks whether more
observability can be bought without new hardware; \cref{ch:overcoming} turns each declared
limitation into its own experiment; \cref{ch:breakthrough} is the project's furthest reach,
the ground-truth-validated certificate; and \cref{ch:economics} asks, cautiously, what any of
it is worth. \Cref{ch:conclusions} closes. The appendices hold the decision register (every
choice and its justification), the dated milestone log, and a plain-language deep dive on how
each machine-learning method works (\cref{ch:ml-appendix}).
